\chapter{Introduction}

\justify

Magna International is one of the world's largest automotive component suppliers, producing structural and powertrain parts via High-Pressure Die Casting (HPDC). In HPDC, molten aluminium is injected into reusable steel dies at extreme pressures, producing complex near-net-shape components at high production rates. The two primary defect classes targeted in this project are Cracks --- linear fractures posing structural safety risks --- and Cold Shuts --- flow lines where two metal fronts fail to fuse. Automated vision-based detection is critical for reliable, consistent quality assurance at production speed.

\section{Preamble (Provide Introduction of the project)}
\justify

This project develops an end-to-end automated surface defect detection system for HPDC aluminium castings in partnership with Magna International. The system uses RF-DETR Medium, a Detection Transformer with a DINOv2 ViT-B backbone, fine-tuned on a patch-based dataset derived from 256 high-resolution HPDC casting images. The core innovation is a sliding-window patch pipeline that converts full 4096$\times$3000 pixel images into 1024$\times$1024 training patches, making it possible to detect cracks that occupy less than 0.1\% of the image area --- a scale at which standard full-image detectors fail entirely.

The project is structured in two phases. Phase 1 (current) covers bounding box detection and localisation of Crack and Cold Shut defects. Phase 2 (deferred) will extend to pixel-level segmentation and physical crack length measurement in millimetres.

\section{Motivation}
\justify

The motivation for this project arises from a combination of industrial necessity and technical challenge:

\begin{itemize}[leftmargin=2em]
    \item \textbf{Safety criticality:} HPDC components are used in structural automotive applications. Undetected cracks can propagate under mechanical and thermal loads, leading to part failure and safety hazards downstream.
    \item \textbf{Production scale:} A single experienced inspector evaluating thousands of parts per shift introduces fatigue-induced inconsistency. Automated detection provides reproducible results at machine speed.
    \item \textbf{Extreme scale challenge:} Exploratory data analysis revealed that crack bounding boxes occupy as little as 0.065\% of the total image area (roughly $90\times90$ pixels on a 12.3~MP image). This makes crack detection one of the hardest small-object detection problems in industrial computer vision.
    \item \textbf{Commercial licensing:} YOLO-family models (versions 8 and beyond) use the AGPL-3.0 licence, which requires derivative commercial products to also be open-sourced --- incompatible with industrial deployment. RF-DETR (Apache 2.0) is the commercially viable alternative.
\end{itemize}

\section{Literature Study}
\justify

Automated defect detection in manufacturing has advanced substantially with deep convolutional neural networks. Traditional approaches relied on hand-crafted features --- HOG descriptors, LBP histograms, Gabor filters --- combined with classical classifiers such as SVM or AdaBoost \cite{xu2020}. Deep learning methods have supplanted these by learning hierarchical representations directly from labelled data.

The Detection Transformer (DETR) paradigm \cite{carion2020} reformulated object detection as a direct set prediction problem, eliminating hand-designed anchor boxes and non-maximum suppression. Deformable DETR \cite{zhu2020} improved convergence and small-object performance through sparse deformable attention. \textbf{RF-DETR} \cite{robinson2025} builds on this lineage with a DINOv2 ViT-B backbone \cite{oquab2023}, multi-scale feature projector, and deformable decoder, making it ideal for fine-tuning on small industrial datasets with challenging object scales. Xu et al.\ \cite{xu2020} demonstrated the viability of machine learning for internal casting crack detection, motivating this surface-defect extension.

\section{Problem Definition}
\justify

Given 256 high-resolution HPDC casting images (4096$\times$3000 pixels) annotated with COCO polygon-level bounding boxes for two defect classes (Crack and Cold Shut), the problem is defined as:

\begin{itemize}[leftmargin=2em]
    \item \textbf{Primary:} Train an object detection model that accurately detects and localises Crack and Cold Shut defects, measured by COCO mAP50-95 on a held-out test set.
    \item \textbf{Key constraint:} Crack bounding boxes occupy $<0.1\%$ of image area, making full-image detection infeasible. A patch-based sliding-window strategy is required.
    \item \textbf{Licence constraint:} The model must use a licence compatible with commercial deployment (Apache 2.0 preferred). YOLO-family models are excluded.
    \item \textbf{Phase 1 scope:} Bounding box detection only. Segmentation and crack length measurement are deferred to Phase 2.
\end{itemize}

\section{Objectives of the project}
\justify

\begin{enumerate}[leftmargin=2em]
    \item Detect and classify surface defects (Crack and Cold Shut) in HPDC casting images using a commercially deployable deep learning model.
    \item Produce accurate bounding box localisations suitable for downstream quality control actions.
    \item Establish a reproducible patch-based training pipeline for extreme small-object detection.
    \item Quantify the impact of dataset engineering choices (patch size, overlap, balancing, augmentation) on detection performance through structured experiments.
    \item Validate RF-DETR as an Apache 2.0 licensed alternative to YOLO, achieving EMA mAP50-95 $\geq$ 0.20 on the test set.
\end{enumerate}
